\chapter{钱与度量}\label{ch:g2:measure}

尺子、厨房秤、量杯、硬币、时钟和日历：
本章讲的是日常生活中的数字——
有多长、多重、能装多少、要花多少钱，以及
现在是几点。

\section{长度、质量、容积}

\begin{definition}[日常单位]\label{def:g2:measure:units}
\begin{itemize}
  \item \emph{长度}：小物体用厘米（cm），大物体用
        米（m），其中 $1$ m $= 100$ cm；
  \item \emph{质量}\index{质量}：轻的物体用克（g），重的
        物体用千克（kg），且 $1$ kg $= 1000$ g；
  \item \emph{容积}\index{容积}（容器能装多少）：升（L）。
\end{itemize}
更小的单位及其精确规则在\cref{ch:g3:measure}中介绍。
\end{definition}

\begin{example}[用哪个单位？]\label{ex:g2:measure:choose}
蜡笔约 $9$ cm。门约 $2$ m。苹果约 $150$
g。一大袋面粉重 $1$ kg。一瓶牛奶的\omterm{def:g2:measure:units}{容积}为 $1$ L。先猜单位，
再测量——这样可以避免出现“门高 $2$
cm”这种荒唐答案。
\end{example}

\section{钱}

\begin{method}[凑出金额]\label{met:g2:measure:money}
用面值为 $1$、$2$、$5$、$10$、$20$ 的硬币和纸币：
\begin{enumerate}
  \item 先用能放下的最大面值，再补足剩下的金额；
  \item 凑出 $27$：一个 $20$、一个 $5$ 和一个 $2$；
  \item 凑法往往不止一种——$27$ 也可以是 $10 + 10 + 5 +
        2$。
\end{enumerate}
\end{method}

\begin{omfigure}
\begin{tikzpicture}[scale=0.9]
  \draw[thick, fill=omProp!20] (0,0) rectangle (1.3,0.7);
  \node[font=\small] at (0.65,0.35) {$20$};
  \node[font=\large] at (1.9,0.35) {$+$};
  \draw[thick, fill=omMeth!20] (2.85,0.35) circle (0.35);
  \node[font=\small] at (2.85,0.35) {$5$};
  \node[font=\large] at (3.8,0.35) {$+$};
  \draw[thick, fill=omDef!20] (4.7,0.35) circle (0.3);
  \node[font=\small] at (4.7,0.35) {$2$};
  \node[font=\large] at (5.6,0.35) {$=$};
  \node[font=\normalsize] at (6.5,0.35) {$27$};
\end{tikzpicture}

{\small 用一张纸币和两枚硬币凑出 $27$——另外还有其他凑法，
例如 $10 + 10 + 5 + 2$。}
\end{omfigure}

\begin{example}[找零]\label{ex:g2:measure:change}
一本书要 $14$；你用 $20$ 付钱，找回的钱是 $20 - 14 =
6$，用“往上数”的方法算出：$14 \to 20$ 相差 $6$
（见\cref{met:g2:subtraction:countup}）。
\end{example}

\section{时间}

\begin{method}[半点和一刻]\label{met:g2:measure:clock}
长针表示的是整点过了几分：
\begin{enumerate}
  \item 竖直向上（指向 $12$）：整点——$3{:}00$；
  \item 向右转四分之一圈（指向 $3$）：过一刻——$3{:}15$；
  \item 竖直向下（指向 $6$）：半点——$3{:}30$；
  \item 转过四分之三圈（指向 $9$）：差一刻到下一个整点——
        $3{:}45$，即“差一刻四点”。
\end{enumerate}
\end{method}

\begin{omfigure}
\begin{tikzpicture}[scale=0.8]
  \foreach \x/\hh/\mm/\lbl in {0/3/3/{3:15}, 4/3/6/{3:30},
                                8/3/9/{3:45}} {
    \begin{scope}[xshift=\x cm]
    \draw[thick] (0,0) circle (1.2);
    \foreach \h in {1,...,12} {
      \node[font=\tiny] at ({90-\h*30}:0.98) {\h};
    }
    \draw[very thick, omDef] (0,0) -- ({90-(\hh+\mm/12)*30}:0.55);
    \draw[thick, omThm] (0,0) -- ({90-\mm*30}:0.95);
    \node[below, font=\small] at (0,-1.45) {\lbl};
    \end{scope}
  }
\end{tikzpicture}

{\small 过一刻、半点、差一刻：红色长针分别指向
$3$、$6$ 和 $9$。}
\end{omfigure}

\begin{example}[日历]\label{ex:g2:measure:calendar}
一年有 $12$ 个月，依次是：一月、二月、三月、四月，
五月、六月、七月、八月、九月、十月、十一月、十二月。
大多数月份有 $30$ 或 $31$ 天；二月有 $28$（或 $29$）天。
日期要同时说明月和日，例如“三月 $5$ 日”。
\end{example}

\section{练习}

\begin{exercise}[$\star$]\label{exo:g2:measure:1}
选择合适的单位（cm、m、g、kg 或 L）：一支铅笔；
教室的长度；一根羽毛；一个书包；一个水瓶。
\end{exercise}

\begin{exercise}[$\star$]\label{exo:g2:measure:2}
换算：$2$ m 化为 cm；$300$ cm 化为 m；$1$ kg 化为 g。
\end{exercise}

\begin{exercise}[$\star$]\label{exo:g2:measure:3}
用硬币和纸币凑出下列金额：$16$；\quad $23$；\quad
$35$。每个金额各给出一种凑法。
\end{exercise}

\begin{exercise}[$\star$]\label{exo:g2:measure:4}
计算找零：用 $10$ 元买一个 $6$ 元的玩具；用 $20$ 元买一本 $17$ 元的
图书；用 $50$ 元买一个 $32$ 元的游戏。
\end{exercise}

\begin{exercise}[$\star$]\label{exo:g2:measure:5}
读出时间：长针指向 $6$，短针在 $4$ 和 $5$ 之间；
长针指向 $3$，短针刚过 $8$。
\end{exercise}

\begin{exercise}[$\star$]\label{exo:g2:measure:6}
画出表示两点半的钟面。两根指针分别指向哪里？
\end{exercise}

\begin{exercise}[$\star$]\label{exo:g2:measure:7}
说出下列月份：五月之后的一个月；一月之前的一个月；你
生日所在的月份。
\end{exercise}

\begin{exercise}[$\star$]\label{exo:g2:measure:8}
莉亚买一杯饮料花了 $3$ 元，买一块蛋糕花了 $4$ 元。她用一张 $10$ 元的
纸币付款。应找回多少钱？
\end{exercise}

\begin{exercise}[$\star$]\label{exo:g2:measure:9}
哪个更重：一袋 $1$ kg 的糖，还是 $800$ g 的苹果？相差多少
克？
\end{exercise}

\begin{exercise}[$\star\star$]\label{exo:g2:measure:10}
用恰好四枚硬币或纸币凑出 $40$（从 $1$、$2$、
$5$、$10$、$20$ 中选取）。求出两种不同的凑法。
\end{exercise}
